Die Constraint-basierte Programmierung erlaubt es Entwicklern, erwünschte
Eigenschaften eines Systems zu spezifizieren, und es dem System selbst zu
überlassen, diese aufrechtzuerhalten. Dies resultiert in kompaktem, deklarativem
Quelltext. Dies steht im Gegensatz zu verstreutem Quelltext, welcher bei
imperativer Programmierung oft nötig ist, um Invarianten zu überprüfen und
wiederherzustellen. Constraints bieten in diesem Sinne Flexibilität und
Ausdruckskraft um komplexe Probleme zu lösen und erwünschte Systemzustände zu
erhalten. Trotz seiner Vorteile ist die Constraint-basierte Programmierung
jedoch nicht weit verbreitet, und die imperative Programmierung ist heute die
Norm.

Constraint-basierte Programmiersprachen sowie deren Integration mit
Allzweckprogrammiersprachen, insbesondere aus dem imperativen Paradigma, haben
eine lange Forschungsgeschichte. Allerdings beschränkt sich die Integration
bisheriger Arbeiten zumeist auf reine Interoperabilität der Paradigmen, statt
deren Abstraktionsmechanismen zu vereinheitlichen. Das führt dazu, das
Constraint-Quelltext und imperativer Quelltext vermischt werden, aber nicht
einheitlich benutzt und entwickelt werden können. Das schränkt die
Wiederverwendbarkeit von Modulen ein, da diese explizit für das eine oder andere
Paradigma optimiert werden müssen, oder redundante Definitionen erfordern, um in
beiden Paradigmen einsetzbar zu sein. Desweiteren zwingt der scharfe Bruch beim
Paradigmenwechsel die Programmierer dazu, beide Paradigmen voll zu verstehen,
bevor sie sinnvoll eingesetzt werden können.

\begin{sloppypar}
  In unserer Arbeit haben wir ein Design für eine Sprachfamilie von
  Objekt-Constraint-basierten Sprachen entworfen, welches wir \emph{Babelsberg}
  nennen. In unserem Design sind die Konstrukte zur Abstraktion vereinheitlicht,
  indem nur objektorientierte Abstraktionen und Methodendefinitionen für sowohl
  imperativen als auch deklarativen Quelltext verwendet werden. Genau wie
  imperativer Fehlerprüfcode können Constraints nun objektorientierte
  Sprachkonstrukte verwenden, und das System kann diese in eine deklarative Form
  übersetzen und für den Rest der Laufzeit kontinuierlich prüfen und
  sicherstellen. Zu diesem Design präsentieren wir außerdem eine Semantik,
  welche Sprachentwicklern als detaillierte Grundlage zur Integration unseres
  Designs mit existierenden, objektorientierten Sprachen dienen kann. Die
  Korrektheit dieser Semantik haben wir für Schlüsseltheoreme theoretisch
  bewiesen, sowie praktisch mithilfe einer ausführbaren Semantik
  gezeigt. Mithilfe dreier konkreter Implementierungen unseres Designs zeigen
  wir weiterhin dessen Praktikabilität. Um Entwicklern aus dem imperativen
  Paradigma den Einsatz von Constraints zu ermöglichen, ohne dass diese die
  spezifischen Limitierungen der verschiedenen automatischen
  Constraint-Lösungsstrategien verstehen müssen, präsentieren wir außerdem eine
  Architektur die es ermöglicht, automatisch verschiedene Lösungsstrategien zu
  kombinieren, um auch komplexe Probleme zu lösen. Schlussendlich diskutieren
  wir anhand der konkreten Implementierungen deren Ausführungsgeschwindigkeit
  unter verschiedenen Bedingungen, sowie Erfahrungen in der Entwicklung und
  Benutzung von Objekt-Constraint-basierten Sprachen.
\end{sloppypar}

Wir glauben, dass unser Ansatz einen nützlichen Schritt darstellt, um das
automatische Lösen von Constraints zu einem nützlichen Werkzeug für imperativen
Programmierer zu machen. Das zeigen wir auch anhand von Beispielanwendungen, in
denen wir unsere Implementierungen des Babelsberg Designs in einer Reihe von
Anwendungsdomänen wie interaktive Grafiken, physikalische Simulationen,
Datenübertragung mit Constraints über die Performance, sowie interaktive
Puzzlespiele.

%%% Local Variables:
%%% mode: latex
%%% TeX-master: "Felgentreff_2015_Babelsberg"
%%% End:
